\documentclass[10pt,a4paper]{amsart}
\usepackage{amsmath,amssymb,amsthm,graphicx}
\usepackage{verbatim}
\newenvironment{reference}{\comment}{\endcomment}
%\newenvironment{reference}{}{}
\newenvironment{slogan}{\comment}{\endcomment}
\newenvironment{history}{\comment}{\endcomment}

% For commutative diagrams we use Xy-pic
\usepackage[all]{xy}

% We use 2cell for 2-commutative diagrams.
\xyoption{2cell}
% Theorem environments.
%
\theoremstyle{plain}
\newtheorem{theorem}[subsection]{Theorem}
\newtheorem{proposition}[subsection]{Proposition}
\newtheorem{lemma}[subsection]{Lemma}

\theoremstyle{definition}
\newtheorem{definition}[subsection]{Definition}
\newtheorem{example}[subsection]{Example}
\newtheorem{exercise}[subsection]{Exercise}
\newtheorem{situation}[subsection]{Situation}

\theoremstyle{remark}
\newtheorem{remark}[subsection]{Remark}
\newtheorem{remarks}[subsection]{Remarks}

\numberwithin{equation}{subsection}

% Macros
%
\def\lim{\mathop{\mathrm{lim}}\nolimits}
\def\colim{\mathop{\mathrm{colim}}\nolimits}
\def\Spec{\mathop{\mathrm{Spec}}}
\def\Hom{\mathop{\mathrm{Hom}}\nolimits}
\def\Ext{\mathop{\mathrm{Ext}}\nolimits}
\def\SheafHom{\mathop{\mathcal{H}\!\mathit{om}}\nolimits}
\def\SheafExt{\mathop{\mathcal{E}\!\mathit{xt}}\nolimits}
\def\Sch{\mathit{Sch}}
\def\Mor{\mathop{\mathrm{Mor}}\nolimits}
\def\Ob{\mathop{\mathrm{Ob}}\nolimits}
\def\Sh{\mathop{\mathit{Sh}}\nolimits}
\def\NL{\mathop{N\!L}\nolimits}
\def\CH{\mathop{\mathrm{CH}}\nolimits}
\def\proetale{{pro\text{-}\acute{e}tale}}
\def\etale{{\acute{e}tale}}
\def\QCoh{\mathit{QCoh}}
\def\Ker{\mathop{\mathrm{Ker}}}
\def\Im{\mathop{\mathrm{Im}}}
\def\Coker{\mathop{\mathrm{Coker}}}
\def\Coim{\mathop{\mathrm{Coim}}}

% Boxtimes
%
\DeclareMathSymbol{\boxtimes}{\mathbin}{AMSa}{"02}

%
% Macros for moduli stacks/spaces
%
\def\QCohstack{\mathcal{QC}\!\mathit{oh}}
\def\Cohstack{\mathcal{C}\!\mathit{oh}}
\def\Spacesstack{\mathcal{S}\!\mathit{paces}}
\def\Quotfunctor{\mathrm{Quot}}
\def\Hilbfunctor{\mathrm{Hilb}}
\def\Curvesstack{\mathcal{C}\!\mathit{urves}}
\def\Polarizedstack{\mathcal{P}\!\mathit{olarized}}
\def\Complexesstack{\mathcal{C}\!\mathit{omplexes}}
% \Pic is the operator that assigns to X its picard group, usage \Pic(X)
% \Picardstack_{X/B} denotes the Picard stack of X over B
% \Picardfunctor_{X/B} denotes the Picard functor of X over B
\def\Pic{\mathop{\mathrm{Pic}}\nolimits}
\def\Picardstack{\mathcal{P}\!\mathit{ic}}
\def\Picardfunctor{\mathrm{Pic}}
\def\Deformationcategory{\mathcal{D}\!\mathit{ef}}
\usepackage{cleveref}
\crefname{originaldefinition}{original definition}{original definitions}
\Crefname{originaldefinition}{Original definition}{Original definitions}
\makeatletter
\newlabel{definition-lci-gysin}{{0FF3}{source}{}{sourcepointer-1}{}}
\newlabel{definition-lci-gysin@cref}{{[originaldefinition][0][]0FF3}{[1][0][]source}{}{}{}}
\makeatother
\crefname{originallemma}{original lemma}{original lemmas}
\Crefname{originallemma}{Original lemma}{Original lemmas}
\makeatletter
\newlabel{lemma-lci-gysin-easy}{{0FF7}{source}{}{sourcepointer-2}{}}
\newlabel{lemma-lci-gysin-easy@cref}{{[originallemma][0][]0FF7}{[1][0][]source}{}{}{}}
\makeatother
\crefname{originallemma}{original lemma}{original lemmas}
\Crefname{originallemma}{Original lemma}{Original lemmas}
\makeatletter
\newlabel{lemma-identify-chow-for-smooth}{{0FC2}{source}{}{sourcepointer-3}{}}
\newlabel{lemma-identify-chow-for-smooth@cref}{{[originallemma][0][]0FC2}{[1][0][]source}{}{}{}}
\makeatother
\crefname{originallemma}{original lemma}{original lemmas}
\Crefname{originallemma}{Original lemma}{Original lemmas}
\makeatletter
\newlabel{lemma-lci-gysin-commutes}{{0FF6}{source}{}{sourcepointer-4}{}}
\newlabel{lemma-lci-gysin-commutes@cref}{{[originallemma][0][]0FF6}{[1][0][]source}{}{}{}}
\makeatother
\title{The projection formula for a proper morphism of smooth schemes over a field}
\author{The Stacks Project Authors}
\date{stacks-project:chow.tex:local-snapshot-20261012; GNU FDL 1.2 or later}
\begin{document}
\maketitle
\noindent\textit{Editorial note.} All delivered mathematical text is the author's own, copied from the pinned Stacks Project chapter file chow.tex; no statement or proof step was rewritten, reordered or removed.\par\medskip
\noindent\textit{Editorial note.} Preamble provenance: the pinned chapter file opens with an input of the shared preamble and carries no preamble text of its own, so the project's single Stacks-wide preamble staged at sources/stacks-chow/preamble.tex is excerpted verbatim as two spans, lines 14-18 (Xy-pic setup) and lines 154-217 (the native theorem declarations and mathematical macros). Nothing was added, removed or redefined; the wrapper class and the amsmath/amssymb/amsthm/graphicx base packages are supplied by the assembler.\par\medskip
\noindent\textit{Editorial note.} The counted statement and proof contain no cross reference, no citation and no included file, so the delivered body is self contained.\par\medskip
\noindent\textit{Editorial note.} Canonical location of the counted proposition: Stacks tag 0FFC (sources/stacks-chow/pinned-original-source.tex line 14482).\par\medskip
\noindent\textit{Allowed standard prerequisites (editorial).}\begin{itemize}
\item Existence of the Gysin map for a local complete intersection morphism. Same chapter, Definition labelled definition-lci-gysin, cited at proof lines 14490-14494. Provides f\textasciicircum{}! for the proper morphism f.
\item Identification of Chow classes of smooth schemes with caps of bivariant classes. Same chapter, Lemma labelled lemma-identify-chow-for-smooth, cited at proof lines 14501-14507 and 14517-14519. Expresses alpha and beta as caps.
\item The value of f\textasciicircum{}! on the fundamental class and its commutation with bivariant classes. Same chapter, Lemma labelled lemma-lci-gysin-easy, cited at proof lines 14499-14501, and Lemma labelled lemma-lci-gysin-commutes, cited at proof lines 14509-14513. Supplies f\textasciicircum{}![Y]\_e = [X]\_d and lets c' be moved past f\textasciicircum{}!.
\item Bivariant classes commute with proper pushforward. Same chapter, the definition of a bivariant class, used at proof lines 14513-14517. Allows c' to be moved past f\_*.
\end{itemize}\medskip
\ifdefined\hypertarget\hypertarget{sourcepointer-1}{}\fi
\noindent\textit{Source reference.} Original definition 0FF3: sources/stacks-chow/pinned-original-source.tex, source lines 13530--13530. Same-chapter pointer into the author environment definition at source line 13530; invoked at proof line 14491 of the counted proof. It is retained only as a navigation pointer; every load-bearing step is inside the delivered span.\par\medskip
\ifdefined\hypertarget\hypertarget{sourcepointer-2}{}\fi
\noindent\textit{Source reference.} Original lemma 0FF7: sources/stacks-chow/pinned-original-source.tex, source lines 13643--13643. Same-chapter pointer into the author environment lemma at source line 13643; invoked at proof line 14496 of the counted proof. It is retained only as a navigation pointer; every load-bearing step is inside the delivered span.\par\medskip
\ifdefined\hypertarget\hypertarget{sourcepointer-3}{}\fi
\noindent\textit{Source reference.} Original lemma 0FC2: sources/stacks-chow/pinned-original-source.tex, source lines 14403--14403. Same-chapter pointer into the author environment lemma at source line 14403; invoked at proof line 14498 of the counted proof. It is retained only as a navigation pointer; every load-bearing step is inside the delivered span.\par\medskip
\ifdefined\hypertarget\hypertarget{sourcepointer-4}{}\fi
\noindent\textit{Source reference.} Original lemma 0FF6: sources/stacks-chow/pinned-original-source.tex, source lines 13613--13613. Same-chapter pointer into the author environment lemma at source line 13613; invoked at proof line 14515 of the counted proof. It is retained only as a navigation pointer; every load-bearing step is inside the delivered span.\par\medskip
\begin{lemma}
\label{lemma-projection-formula}
Let $k$ be a field. Let $f : X \to Y$ be a proper morphism of schemes smooth
over $k$. Then the Gysin map exists for $f$ and
$f_*(\alpha \cdot f^!\beta) = f_*\alpha \cdot \beta$.
\end{lemma}

\begin{proof}
Observe that $X \to X \times_k Y$ is an immersion of $X$ into a scheme
smooth over $Y$. Hence the Gysin map exists for $f$
(Definition \ref{definition-lci-gysin}).
To prove the formula we may decompose $X$ and $Y$ into their
connected components, hence we may assume $X$ is smooth over $k$
and equidimensional of dimension $d$ and $Y$ is smooth over $k$
and equidimensional of dimension $e$. Observe that
$f^![Y]_e = [X]_d$ (see for example Lemma \ref{lemma-lci-gysin-easy}).
Write $\alpha = c \cap [X]_d$ and $\beta = c' \cap [Y]_e$,
see Lemma \ref{lemma-identify-chow-for-smooth}. We have
\begin{align*}
f_*(\alpha \cdot f^!\beta)
& =
f_*(c \cap f^!(c' \cap [Y]_e)) \\
& =
f_*(c \cap c' \cap f^![Y]_e) \\
& =
f_*(c \cap c' \cap [X]_d) \\
& =
f_*(c' \cap c \cap [X]_d) \\
& =
c' \cap f_*(c \cap [X]_d) \\
& =
\beta \cdot f_*(\alpha)
\end{align*}
The first equality by the result of Lemma \ref{lemma-identify-chow-for-smooth}
for $X$. By Lemma \ref{lemma-lci-gysin-commutes} we know that $f^!$
commutes with $c'$. The commutativity of the intersection
product justifies switching the order of capping $[X]_d$ with $c$ and $c'$
(via the lemma). Commuting $c'$ with $f_*$ is allowed as $c'$
is a bivariant class. The final equality is again the lemma.
\end{proof}

\end{document}
